\part{Graphics and comparisons}

\section{Signed bars, targets, intervals and uncertainty}
\begin{demo}{A signed point estimate and asymmetric uncertainty}
\gradbar[min=-50,max=100,width=65mm,
  target=40,band={20,60},error minus=8,error plus=12]{35}

\gradbar[min=-50,max=100,width=65mm,shape=lollipop]{-25}

\gradrange[min=-50,max=100,width=65mm,range point=10]{-20}{35}
\end{demo}
Error magnitudes are nonnegative distances from the observation, not absolute
endpoints. Explain whether they represent standard deviation, standard error or
confidence intervals in the table caption. Targets and bands are references, not
automatically estimated uncertainty. A band must have ordered endpoints within
the scale; invalid reference bands and targets cause errors.

Floating intervals require lower$\le$upper; equal endpoints draw a cap, not an
inflated bar. Both endpoints must be present or both missing. A range point must
lie inside the original endpoints. Intervals reject error whiskers and quality
thresholds. Default labels show the supplied endpoints.
Lollipops show a point for zero, but no point for missing input.

Without an explicit theme, negative bars retain the solid lightred default.
Explicit themes select complementary negative colors. To override the negative gradient,
set its two endpoints explicitly:
\begin{demo}{Independent negative gradient}
\gradbar[min=-100,max=100,width=50mm,
  left color=teal,right color=teal!25,
  negative fill=gradient,
  negative left color=magenta!25,
  negative right color=magenta]{-60}
\end{demo}
The negative left endpoint is the outer side of the bar and the right endpoint is
the side next to zero. `negative color` remains the fallback solid color; threshold
and print modes continue to override the gradient.


\subsection{Two signed observations on one track}
\verb|\gradsignedbar[options]{negative}{positive}| is built into core,
including \key{libraries=none}. Its arguments are independent observations, not
a net total or a computed difference. The first must be nonpositive and the
second nonnegative. Set $\mathrm{min}<0<\mathrm{max}$; asymmetric and named
shared scales are supported.
\begin{demo}{One track with two observations and uncertainties}
\gradsignedbar[min=-50,max=50,width=50mm,
  error=8,theme=blue]{-24}{24}

\gradsignedbar[min=-20,max=80,width=50mm,
  error minus=2,error plus=5,theme=teal]{-10}{40}

\gradsignedbar[min=-50,max=50,width=50mm,
  theme=accessible]{NA}{24}
\end{demo}
The track, zero line and reference marks are drawn once. Error distances apply to
each observation separately; intervals may cross zero and retain the ordinary
\key{overflow} policy. Either argument may independently be \key{NA} or empty;
zero means zero length, not missing. The default outside label shows
``negative / positive'', formatting each value and unit independently. Supported
labels are \key{outside}, \key{left} and \key{none}; \key{text} replaces the whole
label. This command accepts ordinary bars only and adds no native table column
type. Its automatic description includes both observations; an enclosing TikZ
picture owns the description when embedded. Custom scalar backends need explicit
\key{backend fallback=tikz}.
